\chapter{ニューロンは部品箱}
% full version: chapters/18_eetypes.tex

第1章では、ニューロンを放出する物質で仕分けた。エンジニアなら別の仕分け方をするだろう。それぞれがどんな素子として働いているか、である。すると、同じ見た目の下に電子部品の箱がまるごと隠れていることがわかる。

\section*{電子部品の箱}

第2章の穴のあいたバケツは、漏れのある積分器である。短い窓の中で信号を足し合わせる。窓をごく小さくすると、同時検出器になる。変化にだけ応答してすぐ黙るニューロンは、変化を通すフィルタだ。自分で動くニューロンはメトロノーム、つまり発振器であり、入力がそのテンポを調整するなら電圧制御発振器である。カチッのバーストを出すニューロンは、「ツー」を書くバースト発生器だ。音波に合わせてカチッと鳴るニューロンは位相検出器である。

\pencilfig{18}{\pencilart{18}}{同じ見た目の下に電子部品の箱がまるごと隠れている。抵抗、コンデンサ、コイル、発振器。}

これまで話していない素子もある。\emph{共振器}はブランコに似ている。ちょうどよいタイミングで押したときだけ大きく揺れるブランコのように、決まったリズムの押しに最もよく応える。ニューロンの内部では、変化に逆らうゆっくりした電流がばねの役を果たす。\emph{漏れのない積分器}は、漏れのないコンデンサのように値を長く保つ。あなたが横を見ているときに目の位置を保つネットワークは、こうして働いている。ただしそのためにはフィードバックをほぼ完璧に合わせなければならず、これは高くつく調整である。\emph{ラッチ}は、短い一押しで長く働き続け、オフにされるまで止まらないニューロンだ。ロックつきのスイッチのようなものである。筋肉を制御するニューロンでは、セロトニンがこのモードをオンにする。

\section*{カメレオン素子}

この仕分けの最大の発見は、これらが別々の部品ではなく、同じ1つの部品の別々の\emph{モード}だということだ。同じニューロンが、放送局に何を言われるかしだいで、昼は積分器、夜はバースト発生器になりうる。エンジニアなら脳を書き換え可能な論理アレイと呼ぶだろう。回路は1つで、部品の役割は設定が変える。

ニューロンがこうしたすべてのモードで働けることは測定されている。脳がモードの切り替えを計算にどう使っているかは、大部分が仮説である。

\fullversion{18}
